\section{Metodologia}
Esta seção descreve a metodologia adotada para o desenvolvimento da ferramenta \offscan, detalhando o ambiente de desenvolvimento, a arquitetura do sistema, as funcionalidades implementadas, os procedimentos de teste e os métodos de coleta de dados.

\subsection{Ambiente de Desenvolvimento}
O desenvolvimento do \offscan foi realizado em ambiente Linux (distribuições Debian 12). Este sistema operacional foi selecionado devido ao suporte nativo a \textit{sockets} brutos (\textit{raw sockets}) e chamadas de sistema (\textit{syscalls}) de baixo nível, requisitos mandatórios para operações ofensivas de rede. 

A linguagem de programação Go foi selecionada para o desenvolvimento desta ferramenta devido a características técnicas que se alinham aos requisitos de sistemas de segurança ofensiva. A simplicidade de sua sintaxe reduz a complexidade do código, facilitando a manutenção e a depuração do software. Por ser uma linguagem compilada diretamente para código de máquina, Go apresenta desempenho em tempo de execução adequado para operações de alta velocidade, como varreduras de diversas redes simultaneamente. Adicionalmente, o modelo de concorrência baseado em Goroutines viabiliza a execução eficiente de tarefas paralelas, enquanto seu ecossistema nativo oferece suporte à manipulação de protocolos de redes de computadores. Desse modo, os recursos da linguagem suprem as demandas de desempenho, concorrência e portabilidade necessárias para a arquitetura do \offscan.

As principais ferramentas e bibliotecas utilizadas incluem:
\begin{itemize}
    \item \textbf{Adaptador de Rede Sem Fio}: Dispositivo de hardware equipado com \textit{chipset} homologado para suporte nativo ao modo monitor (\textit{monitor mode}) e injeção de pacotes (\textit{packet injection}), condições indispensáveis para a interceptação de quadros de gerenciamento e execução de ataques na camada de enlace.

    \item \textbf{Sistema Operacional Linux}: Ambiente de execução obrigatório para o funcionamento da ferramenta, uma vez que a implementação depende de chamadas de sistema (\textit{syscalls}) específicas do kernel Linux para captura de pacotes, manipulação de \textit{sockets} e alteração das configurações da interface de rede sem fio (como a ativação do modo monitor).
    
    \item \textbf{Compilador Go}: Sistema de compilação e gerenciador de pacotes do ecossistema Go, utilizando a versão 1.25.0 ou superior.
    
    \item \textbf{golang.org/x/sys}: Repositório oficial do ecossistema Go (externo à biblioteca padrão), empregado para realizar chamadas de sistema diretas e configurar os \textit{sockets}.
\end{itemize}

\subsection{Arquitetura da Ferramenta}
A arquitetura do \offscan segue um padrão de projeto modular, organizado em componentes especializados que colaboram de forma coordenada para executar operações ofensivas. Seus principais módulos consistem em:

\begin{itemize}
    \item \textbf{Motores (\textit{Engines})}: Blocos lógicos que coordenam e executam as operações fim da ferramenta, utilizando as funções dos demais módulos para consolidar rotinas de varredura ou ataque.
    
    \item \textbf{Construtores de Pacotes/Quadros}: Componentes responsáveis pela montagem manual de pacotes de rede em múltiplas camadas (ARP, IP, TCP, ICMP) e de quadros do padrão IEEE 802.11.
    
    \item \textbf{Sockets Brutos}: Módulo encarregado do envio direto dos pacotes e quadros injetados na interface de rede.
    
    \item \textbf{Sniffer}: Componente encarregado da captura e interceptação de pacotes ou quadros em tempo real diretamente de uma interface de rede.
    
    \item \textbf{Dissecadores de Pacotes/Quadros}: Funções responsáveis pela decodificação e extração de campos e metadados dos pacotes e quadros capturados.
    
    \item \textbf{Geradores}: Blocos lógicos encarregados de produzir iteradores e valores pseudoaleatórios (como portas de origem ou identificadores).
    
    \item \textbf{Provedor de Informações}: Subsistema encarregado de coletar metadados das interfaces físicas de rede e parâmetros operacionais do sistema subjacente.
    
    \item \textbf{Modelos (\textit{Models})}: Componente que centraliza a definição de estruturas (\textit{structs}) e tipos customizados do sistema, atuando na padronização de dados, facilitação da legibilidade e garantia de consistência durante o desenvolvimento.
    
    \item \textbf{Utilitários}: Funções auxiliares isoladas que fornecem rotinas genéricas e reutilizáveis em múltiplos módulos do sistema.
\end{itemize}

\subsection{Fluxo de Execução}
A ferramenta opera como uma aplicação de interface de linha de comando (CLI), na qual cada funcionalidade é invocada por meio de subcomandos específicos. O fluxo geral compreende quatro etapas sequenciais: (1) análise sintática (\textit{parsing}) dos argumentos e parâmetros de entrada; (2) inicialização e acoplamento dos módulos exigidos; (3) execução do motor principal correspondente à operação requisitada (varredura, ataque ou análise); e (4) exibição estruturada dos resultados ao usuário.
